\section{Excluding degenerate optimizers}\label{sec:degenerate-optimizers}

In view of \Cref{prop:nondegenerate-multipodality}, the following
proposition completes the proof of \Cref{thm:fixed-densities}.

\begin{prop}\label{prop:optimizer-nondegeneracy}
Let $H$ be a graph with a vertex of degree at least two, let $\eps\in(0,1)$,
and let $\eps^m<\tau<\tau_{\max}(\eps)$. Every minimizer
of~\eqref{eq:fixed-problem} satisfies the constraint qualification.
\end{prop}

We first exclude degeneracy for almost every $\tau$ in
\Cref{cor:almost-every-tau}. To treat an arbitrary density $\tau$,
we choose optimizers with an ordered structure at densities approaching
$\tau$ from below and extract a convergent subsequence. The limit is
again an optimizer with this structure, although its multipodality is
not yet established. \Cref{lem:ordered-nondegenerate} shows that every
optimizer with an ordered structure is nondegenerate, so the limit satisfies the
constraint qualification. Without uniqueness, however, other optimizers
at $\tau$ could still be degenerate. In
\Cref{sec:removing-exceptional-densities}, we therefore use the perturbation
directions at the limit to bound the increase in $\Psi_\eps$ and rule out
this possibility.

\subsection{Multipodality at almost every density}\label{sec:almost-every-multipodality}

\begin{lemma}\label{lem:degenerate-forms}
Let $W$ be a degenerate minimizer of~\eqref{eq:fixed-problem} and put
$\Omega:=\{(x,y):0<W(x,y)<1\}$.  Then $\nu(\Omega)>0$, and there is a constant $\kappa$
with $\mathcal D_{H,W}=\kappa$ almost everywhere on $\Omega$.
\end{lemma}

\begin{proof}
Let $W_{\max}$ attain $\tau_{\max}(\eps)$.  By continuity of the
$H$-density along the segment from $\eps$ to $W_{\max}$, there is
$\theta\in(0,1)$ such that
\[
  U:=\eps+\theta(W_{\max}-\eps),
  \qquad e(U)=\eps,\quad t(H,U)=\tau.
\]
The values of $U$ lie in a compact subinterval of $(0,1)$,
so $s(U)>0$.
Optimality and~\eqref{eq:entropy-identity} give $s(W)\ge s(U)>0$, which
forces $\nu(\Omega)>0$.

For every $\gamma>0$ such that $\Omega_\gamma(W)$ has positive measure, the
characterization following \Cref{def:constraint-qualification} shows
that $\mathcal D_{H,W}$ is essentially constant on $\Omega_\gamma(W)$.  These
constants agree, because the sets $\Omega_\gamma(W)$ increase as $\gamma$
decreases and have positive measure.  Since
$\Omega=\bigcup_{\gamma>0}\Omega_\gamma(W)$, the common value $\kappa$ satisfies
$\mathcal D_{H,W}=\kappa$ almost everywhere on $\Omega$.
\end{proof}

\begin{lemma}\label{lem:degenerate-smoothing}
Let $W$ be a degenerate minimizer of~\eqref{eq:fixed-problem} at $\tau$.
Let $W_s:=(1-s)W+s\eps$, and
denote its cost saving and density loss by
\[
  g_s:=I_\eps(W)-I_\eps(W_s),
  \qquad
  d_s:=\tau-t(H,W_s)
  \qquad(0\le s\le1).
\]
Then $g_s$ and $d_s$ are positive for $0<s<1$.
Moreover, $d_s$ is a polynomial in $s$ and
\[
  d_s=a_k s^k+O(s^{k+1}),\qquad
  \frac{g_s}{d_s}\longrightarrow\infty
  \quad\text{as }s\downarrow0
\]
for some $k\ge1$ and $a_k>0$.
\end{lemma}

\begin{proof}
We have $e(W_s)=\eps$, and $W\not\equiv\eps$ because $\tau>\eps^m$, so
$I_\eps(W)>0$.  Convexity of $J_\eps$ and $J_\eps(\eps)=0$ give
\[
  I_\eps(W_s)\le(1-s)I_\eps(W),
  \qquad g_s\ge s I_\eps(W)>0.
\]
By the comparison in \Cref{lem:fixed-monotonicity}, a graphon with edge
density $\eps$ and cost below $I_\eps(W)$ cannot have $H$-density at least
$\tau$; hence $d_s>0$.  Since $s\mapsto d_s$ is a polynomial
vanishing at $s=0$ and positive on $(0,1)$, its first nonzero
coefficient $a_k$ satisfies $k\ge1$ and $a_k>0$.

If $\{W=0\}\cup\{W=1\}$ has positive measure, the endpoint values give an
additional logarithmic saving.  For small $s>0$,
\[
  J_\eps(0)-J_\eps(s\eps) =\eps s\log(1/s)+O(s),\qquad J_\eps(1)-J_\eps(1-s(1-\eps)) =(1-\eps)s\log(1/s)+O(s).
\]
Since interpolation toward $\eps$ decreases $J_\eps$ at every other
point as well, integration gives
\[
  g_s\ge c s\log(1/s),\qquad d_s\le C s,
\]
after decreasing $c>0$ and increasing $C>0$ if necessary.
Thus $g_s/d_s\to\infty$ as $s\downarrow0$.

Otherwise $0<W<1$ almost everywhere, so $\mathcal D_{H,W}=\kappa$ almost everywhere by
\Cref{lem:degenerate-forms} and
\[
  \frac{\dd}{\dd s}t(H,W_s)\Big|_{s=0}
  =\int\mathcal D_{H,W}\,(\eps-W)=\kappa\bigl(\eps-e(W)\bigr)=0.
\]
Thus the first order term of $d_s$ vanishes, which implies $k\ge2$,
and the bounds
\[
  g_s\ge s I_\eps(W), \qquad
  d_s\le C s^2
\]
again imply $g_s/d_s\to\infty$ as $s\downarrow0$.
\end{proof}

The following lemma turns the comparison in \Cref{lem:degenerate-smoothing}
into an infinite left slope of $\Psi_\eps$.

\begin{lemma}\label{lem:degenerate-slope}
If~\eqref{eq:fixed-problem} has a degenerate minimizer at
$\tau\in(\eps^m,\tau_{\max}(\eps))$, then
\[
  \lim_{\tau'\uparrow\tau}\frac{\Psi_\eps(\tau)-\Psi_\eps(\tau')}{\tau-\tau'}=\infty.
\]
\end{lemma}

\begin{proof}
Keep the notation of \Cref{lem:degenerate-smoothing}.  The polynomial
density loss has positive leading coefficient, so
\[
  \frac{\dd}{\dd s}d_s=ka_k s^{k-1}+O(s^k)>0
\]
for small $s>0$.  Every $\tau'<\tau$ sufficiently close to $\tau$
therefore has a unique small $s>0$ with $d_s=\tau-\tau'$.
The graphon $W_s$ is feasible for~\eqref{eq:fixed-problem} at $\tau'$,
thus $\Psi_\eps(\tau')\le I_\eps(W_s)$, and
\[
  \frac{\Psi_\eps(\tau)-\Psi_\eps(\tau')}{\tau-\tau'}
  \ge\frac{I_\eps(W)-I_\eps(W_s)}{d_s}
  =\frac{g_s}{d_s}.
\]
As $\tau'\uparrow\tau$ we have $s\downarrow0$, and the right-hand side
tends to infinity by \Cref{lem:degenerate-smoothing}.
\end{proof}

\begin{cor}\label{cor:almost-every-tau}
The set of $\tau\in(\eps^m,\tau_{\max}(\eps))$ for
which~\eqref{eq:fixed-problem} has a degenerate minimizer has Lebesgue measure
zero.  Consequently, for almost every such $\tau$, every maximizer
of~\eqref{eq:entropy-maximization} is multipodal with ordered blocks.
\end{cor}

\begin{proof}
By \Cref{lem:fixed-monotonicity}, $\Psi_\eps$ is increasing, so Lebesgue's
theorem implies that it is
differentiable with a finite derivative at almost every point of
$(\eps^m,\tau_{\max}(\eps))$.
At such a point, the left difference
quotients converge to the finite derivative, so by
\Cref{lem:degenerate-slope} no degenerate minimizer exists there.
At every such point, \Cref{prop:nondegenerate-multipodality} applies to every maximizer.
\end{proof}

\subsection{Ordered optimizers are nondegenerate}\label{sec:ordered-nondegeneracy}

Without a uniform bound on the number of blocks, compactness does not
imply that a limit of the minimizers in \Cref{cor:almost-every-tau} is
multipodal. We therefore prove nondegeneracy for any minimizer with
an ordered structure, as defined in \Cref{def:ordered-structure}.

\begin{lemma}[Global row maximality]\label{lem:degenerate-insertion}
Let $W$ be a degenerate minimizer of~\eqref{eq:fixed-problem} at
$\tau\in(\eps^m,\tau_{\max}(\eps))$, and let $\kappa$ be as in
\Cref{lem:degenerate-forms}.  Then
\begin{equation}\label{eq:degenerate-insertion}
  \mathcal B_W(u)
  :=a_H(u)-\kappa a_e(u)
  =\mathcal R_{H,W}(u)-N\tau-2\kappa\Bigl(\int_0^1u-\eps\Bigr)\le0
\end{equation}
for every measurable $u:[0,1]\to[0,1]$, with equality for almost every row
$u=W_x$.
\end{lemma}

\begin{proof}
Suppose that $\mathcal B_W(u)>0$ for some measurable $u:[0,1]\to[0,1]$.
Let $W_s$, $g_s$, and $d_s$ be as in \Cref{lem:degenerate-smoothing}.
By \Cref{lem:degenerate-forms} we may fix $\gamma>0$ such that
$\Omega_\gamma(W)$ has positive measure. Choose the direction
\[
  V:=\mathbf1_{\Omega_\gamma(W)},
  \qquad \int V=\nu(\Omega_\gamma(W))>0,
  \qquad \int\mathcal D_{H,W}V=\kappa\int V.
\]

Define $V_{s,\theta,\zeta}$ by~\eqref{eq:vertex-insertion}, using insertion
mass $0\le\theta<1$, row $u$, internal value $\eps$, and old kernel
$W_s+\zeta V$. The parameters $\theta$ and $\zeta$
vary independently near $0$.
The edge density is
\[
  e(V_{s,\theta,\zeta})
  =(1-\theta)^2\Bigl(\eps+\zeta\int V\Bigr)
    +2\theta(1-\theta)\int_0^1u+\theta^2\eps.
\]
For the $H$-density, \eqref{eq:insertion-density-expansion} gives
\[
  t(H,V_{s,\theta,\zeta})
  =(1-\theta)^N t(H,W_s+\zeta V)
    +\theta(1-\theta)^{N-1}\mathcal R_{H,W_s+\zeta V}(u)+O(\theta^2).
\]
Both densities extend to polynomials in $(s,\theta,\zeta)$: each
assignment of $k$ vertices to the new class contributes the mass factor
$\theta^k(1-\theta)^{N-k}$, and the remaining edge factors are $u$, $\eps$, or
$W_s+\zeta V=(1-s)W+s\eps+\zeta V$.
At $s=0$, expanding in $(\theta,\zeta)$ and using
$\int\mathcal D_{H,W}V=\kappa\int V$ gives
\[
  \left.\frac{\partial\mathbf C(V_{s,\theta,\zeta})}{\partial(\theta,\zeta)}
  \right|_{(s,\theta,\zeta)=(0,0,0)}
  =
  \begin{pmatrix}
    a_e(u) & \int V\\[1mm]
    a_H(u) & \kappa\int V
  \end{pmatrix}.
\]
These linear terms show that canceling the change in edge density leaves
\[
  \theta a_e(u)+\zeta\int V=0
  \quad\Longrightarrow\quad
  \theta a_H(u)+\kappa\zeta\int V=\theta\mathcal B_W(u).
\]
Thus $\mathcal B_W(u)$ is the remaining first-order $H$-density change
per unit inserted mass after the perturbation $\zeta V$ of the old graphon
cancels the change in edge density. We use its positivity to restore
the loss $d_s$ with a positive inserted mass.

The Jacobian above has determinant $-(\int V)\mathcal B_W(u)\ne0$.
Applying the implicit function theorem to the polynomial extensions
of the densities therefore gives a local map
$s\mapsto(\theta(s),\zeta(s))$ with $\theta(0)=\zeta(0)=0$ and
\[
  \mathbf C\bigl(V_{s,\theta(s),\zeta(s)}\bigr)
  =\mathbf C(V_{0,0,0})
  =(\eps,\tau).
\]
To bound $\theta(s)$ and $\zeta(s)$, fix $s$ and consider the
density map
\[
  F_s(\theta,\zeta):=\mathbf C(V_{s,\theta,\zeta}),
  \qquad F_s(\theta(s),\zeta(s))=(\eps,\tau),
  \qquad F_s(0,0)=(\eps,\tau-d_s).
\]
The invertibility of the Jacobian above also gives local inverses
$F_s^{-1}$ on a common convex neighborhood of $(\eps,\tau)$ for
small $s$. For $(\eps',\tau')$
in this neighborhood, $F_s^{-1}(\eps',\tau')$ gives the
parameters $(\theta,\zeta)$ solving $F_s(\theta,\zeta)=(\eps',\tau')$.
The derivative of this inverse
depends continuously on $(s,\eps',\tau')$. After shrinking the
neighborhood, there is therefore a constant $C$ such that
\[
  \|D(F_s^{-1})(\eps',\tau')\|_{\mathrm{op}}\le C
\]
uniformly for small $s$ and $(\eps',\tau')$ in the small neighborhood,
using the norm $\|(x,y)\|_1:=|x|+|y|$ on $\RR^2$.  Since
\[
  F_s^{-1}(\eps,\tau)=(\theta(s),\zeta(s)),
  \qquad F_s^{-1}(\eps,\tau-d_s)=(0,0),
\]
the bound on the derivative gives
\[
  |\theta(s)|+|\zeta(s)|
  =\bigl\|F_s^{-1}(\eps,\tau)
           -F_s^{-1}(\eps,\tau-d_s)\bigr\|_1
  \le C\bigl\|(\eps,\tau)-(\eps,\tau-d_s)\bigr\|_1
  =Cd_s.
\]
Expanding the two constraints at $(\theta,\zeta)=(\theta(s),\zeta(s))$ gives
\begin{align*}
  0&=e(V_{s,\theta,\zeta})-e(W_s)
     =a_e(u)\theta(s)+\Bigl(\int V\Bigr)\zeta(s)+O((|\theta(s)|+|\zeta(s)|)^2),\\
  d_s&=t(H,V_{s,\theta,\zeta})-t(H,W_s)
     =a_H(u)\theta(s)+\kappa\Bigl(\int V\Bigr)\zeta(s)+O(d_s^2+s d_s).
\end{align*}
The quadratic remainders are $O(d_s^2)$; replacing the $H$-density
derivatives at $W_s$ by those at $W$ adds $O(s d_s)$.
Both errors are $o(d_s)$.
Subtracting $\kappa$ times the first equation from the second gives
\[
  d_s=\mathcal B_W(u)\theta(s)+o(d_s),
\]
and hence
\[
  \theta(s)=\left(\frac1{\mathcal B_W(u)}+o(1)\right)d_s.
\]
Since $\mathcal B_W(u)>0$ and $d_s>0$ and $d_s\to0$ as $s\downarrow0$,
we have $0<\theta(s)<1$ for small $s>0$.
On the support of $V$, the values
of $W_s$ lie in $[\gamma/2,1-\gamma/2]$ for small $s>0$.
Thus taking $|\zeta(s)|\|V\|_\infty<\gamma/4$ keeps the old
graphon $W_s+\zeta(s)V$ in $[0,1]$.
The inserted row $u$ and the internal value $\eps$ are also admissible.
Thus the constructed graphon $V_{s,\theta(s),\zeta(s)}$ is a feasible graphon.

For the cost, write $\theta=\theta(s)$ and $\zeta=\zeta(s)$. The exact
insertion formula is
\[
  I_\eps(V_{s,\theta,\zeta})
    =(1-\theta)^2I_\eps(W_s+\zeta V)
      +2\theta(1-\theta)I_\eps(u).
\]
Subtracting $I_\eps(W_s)$ and separating the three contributions gives
\[
  I_\eps(V_{s,\theta,\zeta})-I_\eps(W_s)
  =(1-\theta)^2\bigl(I_\eps(W_s+\zeta V)-I_\eps(W_s)\bigr)
  -\theta(2-\theta)I_\eps(W_s)+2\theta(1-\theta)I_\eps(u).
\]
On the support of $V$, both $W_s$ and $W_s+\zeta V$ take
values in $[\gamma/4,1-\gamma/4]$, so define
\[
  M:=\max_{z\in[0,1]}J_\eps(z),\qquad
  L:=\max_{z\in[\gamma/4,1-\gamma/4]}|J_\eps'(z)|.
\]
Outside the support of $V$, the old graphon is unchanged. The mean
value theorem therefore gives
\begin{align*}
  \bigl|I_\eps(W_s+\zeta V)-I_\eps(W_s)\bigr|
  &\le\int_{\Omega_\gamma(W)}
    \bigl|J_\eps(W_s+\zeta V)-J_\eps(W_s)\bigr|\\
  &\le L|\zeta|\int_{\Omega_\gamma(W)}|V|
   =L|\zeta|\|V\|_1.
\end{align*}
Combining these bounds with $0<\theta<1$, we obtain
\begin{align*}
  \bigl|I_\eps(V_{s,\theta,\zeta})-I_\eps(W_s)\bigr|
  &\le(1-\theta)^2L|\zeta|\|V\|_1+\theta(2-\theta)M+2\theta(1-\theta)M\\
  &\le L|\zeta|\|V\|_1+4M\theta
   \le C(|\theta|+|\zeta|)
   = O(d_s).
\end{align*}
Combining this cost bound with the smoothing gain $g_s$ gives
\[
  I_\eps(V_{s,\theta(s),\zeta(s)})
  \le I_\eps(W_s) + Cd_s
  =I_\eps(W)-g_s+ Cd_s,
\]
which is smaller than $I_\eps(W)$ for small $s>0$, because
$g_s/d_s\to\infty$ as $s\downarrow0$ by \Cref{lem:degenerate-smoothing}.
This contradicts the minimality of $W$ and proves~\eqref{eq:degenerate-insertion}.

For existing rows, Fubini's theorem and~\eqref{eq:rooted-average} give
\[
  \int_0^1\mathcal R_{H,W}(W_x)\dd x=N\tau,
  \qquad \int_0^1\int_0^1W_x(y)\dd y\dd x=\eps.
\]
It follows that $\int_0^1\mathcal B_W(W_x)\dd x=0$.
Since $\mathcal B_W(W_x)\le0$, equality holds for almost every row.
\end{proof}

\Cref{lem:degenerate-insertion} makes almost every existing row $u$ of a
degenerate minimizer a global maximizer of $\mathcal B_W$. Keeping $W$
fixed, every admissible pair of rows $u\pm h$ must therefore satisfy
\[
  \mathcal B_W(u+h)+\mathcal B_W(u-h)-2\mathcal B_W(u)\le0.
\]
We will use the ordered structure to find such a pair for which the
left-hand side is strictly positive.

\begin{lemma}\label{lem:ordered-nondegenerate}
Let $W$ be a minimizer of~\eqref{eq:fixed-problem} at
$\tau\in(\eps^m,\tau_{\max}(\eps))$. If $W$ has an ordered structure,
then it satisfies the constraint qualification.
\end{lemma}

\begin{proof}
Suppose that $W$ is degenerate. \Cref{lem:degenerate-forms} gives
$\nu(\{0<W<1\})>0$. By symmetry, the subset where $y\ge x$ also has positive
measure.  Fubini's theorem and \Cref{lem:degenerate-insertion} therefore
give a row $u=W_x$ with $\mathcal B_W(u)=0$ such that the set
\[
  B_x:=\{y \in [0,1]:y\ge x,\ 0<u(y)<1\}
\]
has positive measure.  For $(y,z)\in B_x\times B_x$, monotonicity gives
$W(y,z)\ge W(x,z)=u(z)>0$.
Thus the ordered structure ensures that no edge factor vanishes when
both vertices lie in $B_x$.

To keep both $u+h$ and $u-h$ in $[0,1]$, choose a permissible
change at each coordinate, $h(y):=\min\{u(y),1-u(y)\}$.
Then $h>0$ on $B_x$. Using the definition of $\mathcal B_W$ in
\eqref{eq:degenerate-insertion}, we obtain
\begin{align*}
  \mathcal B_W(u+h)-\mathcal B_W(u)
    &=\mathcal R_{H,W}(u+h)-\mathcal R_{H,W}(u)-2\kappa\int_0^1h,\\
  \mathcal B_W(u-h)-\mathcal B_W(u)
    &=\mathcal R_{H,W}(u-h)-\mathcal R_{H,W}(u)+2\kappa\int_0^1h.
\end{align*}
Adding these identities cancels the linear terms and gives
\begin{equation}\label{eq:rooted-density-difference}
  \mathcal B_W(u+h)+\mathcal B_W(u-h)-2\mathcal B_W(u)
  =\mathcal R_{H,W}(u+h)+\mathcal R_{H,W}(u-h)-2\mathcal R_{H,W}(u).
\end{equation}

We show that the right-hand side of~\eqref{eq:rooted-density-difference}
is strictly positive by expanding the rooted densities.
For each root $v\in V(H)$, the product of row factors satisfy
\[
  \prod_{z\in N_H(v)}(u+h)(x_z)+\prod_{z\in N_H(v)}(u-h)(x_z)
    -2\prod_{z\in N_H(v)}u(x_z)
  =2\sum_{\substack{S\subseteq N_H(v)\\|S|\ge2\text{ even}}}
    \prod_{z\in S}h(x_z)\prod_{z\in N_H(v)\setminus S}u(x_z)\ge0.
\]
Terms with an odd number of factors $h$ cancel, and all remaining terms
are nonnegative. Thus each root contributes a nonnegative amount to
the right-hand side of~\eqref{eq:rooted-density-difference}.
To obtain a positive contribution, choose a vertex $v_0$ with
$d_{v_0}\ge2$ and distinct neighbors $z_1,z_2$.
Keep the term corresponding to $v=v_0$ and $S=\{z_1,z_2\}$, which uses
$h$ on two incident edges. Restricting the labels of all vertices other
than $v_0$ to $B_x$ gives the contribution
\[
  2\int_{B_x^{V(H)\setminus\{v_0\}}}
    h(x_{z_1})h(x_{z_2})
    \prod_{z\in N_H(v_0)\setminus\{z_1,z_2\}}u(x_z)
    \prod_{\substack{ab\in E(H)\\a,b\ne v_0}}W(x_a,x_b)
    \prod_{z\ne v_0}\dd x_z>0,
\]
since $h,u>0$ on $B_x$, $W>0$ on $B_x\times B_x$, and the domain has
measure $\nu(B_x)^{N-1}>0$. Hence
\[
  \mathcal B_W(u+h)+\mathcal B_W(u-h)-2\mathcal B_W(u)>0.
\]
This contradicts the fact that $u$ maximizes $\mathcal B_W$.
\end{proof}

\subsection{Compactness and removal of the exceptional set}\label{sec:removing-exceptional-densities}

\begin{lemma}\label{lem:ordered-compact}
Let $(W_n)$ be graphons with an ordered structure. After
measure-preserving relabelings, some subsequence converges in
$L^1([0,1]^2)$ to a graphon with an ordered structure.
\end{lemma}

\begin{proof}
For each integer $k\ge1$, partition $[0,1]$ into the intervals
$I_a=[(a-1)/k,a/k)$ for $1\le a<k$, and $I_k=[(k-1)/k,1]$.
For a graphon $V$, define its step approximation
$V^{(k)}:[0,1]^2\to[0,1]$ by
\[
  V^{(k)}(x,y):=k^2\int_{I_a\times I_b}V(s,t)\dd s\dd t,
  \qquad (x,y)\in I_a\times I_b.
\]
This is a graphon that is constant on each of the $k^2$ squares.
If $V$ is nondecreasing in each variable, then so is $V^{(k)}$.
On each square, both $V$ and its average lie between the values of $V$
at the lower left and upper right corners. Hence
\[
  \|V-V^{(k)}\|_1
  \le\frac1{k^2}\sum_{a,b=1}^k\Bigl(V\bigl(\tfrac ak,\tfrac bk\bigr)
  -V\bigl(\tfrac{a-1}k,\tfrac{b-1}k\bigr)\Bigr).
\]
Split each corner difference into the increments $V(a/k,b/k)-V((a-1)/k,b/k)$
and $V((a-1)/k,b/k)-V((a-1)/k,(b-1)/k)$. Summing the first-coordinate
increments gives
\[
  \sum_{a,b=1}^k\Bigl(V\bigl(\tfrac ak,\tfrac bk\bigr)
    -V\bigl(\tfrac{a-1}k,\tfrac bk\bigr)\Bigr)
  =\sum_{b=1}^k\Bigl(V\bigl(1,\tfrac bk\bigr)-V\bigl(0,\tfrac bk\bigr)\Bigr)
  \le k.
\]
The sum of the second-coordinate increments is also at most $k$, so
$\|V-V^{(k)}\|_1\le2/k$.

By \Cref{def:ordered-structure}, we may relabel each $W_n$ and modify it
on a null set so that it is symmetric and nondecreasing in both variables.
Keep the notation $W_n$ for these graphons.
For each fixed $k$, the vector of block values of $W_n^{(k)}$ belongs to
the compact set $[0,1]^{k^2}$. We can therefore choose an infinite set
of indices along which all $k^2$ block values converge. Along these
indices, the step graphons converge in $L^1([0,1]^2)$, since their
$L^1$ distance is $k^{-2}$ times the sum
of the absolute differences of their block values.

Choose an infinite index set $\mathcal N_1\subseteq\mathbb N$ along which
$W_n^{(1)}$ converges. Restricting to indices in $\mathcal N_1$, choose an
infinite subset $\mathcal N_2\subseteq\mathcal N_1$ along which $W_n^{(2)}$
converges, and continue. This gives nested infinite index sets
\[
  \mathcal N_1\supseteq\mathcal N_2\supseteq\cdots
\]
such that $W_n^{(k)}$ converges in $L^1([0,1]^2)$ as $n\to\infty$ through
$\mathcal N_k$. Restricting the index set preserves the convergence
obtained at earlier steps. Now choose an increasing sequence $(n_j)_{j=1}^\infty$
with $n_j\in\mathcal N_j$ for each $j$.
For every fixed $k$, we have $n_j\in\mathcal N_j\subseteq\mathcal N_k$
whenever $j\ge k$. Thus $W_{n_j}^{(k)}$ converges for every $k$ along
the same index sequence $(n_j)_{j=1}^\infty$. Keep the notation $(W_n)$ for the
resulting graphon subsequence $(W_{n_j})$. The approximation bound gives
\[
  \|W_n-W_m\|_1
  \le\|W_n-W_n^{(k)}\|_1+\|W_n^{(k)}-W_m^{(k)}\|_1
    +\|W_m^{(k)}-W_m\|_1
  \le\frac4k+\|W_n^{(k)}-W_m^{(k)}\|_1.
\]
Choosing $k$ large and then $n,m$ large shows that $(W_n)$ is Cauchy
in $L^1([0,1]^2)$. By completeness, it has a limit $W_*$ in this space.
The limit is a graphon, since symmetry and the bounds $0\le W_n\le1$
are preserved under $L^1$ convergence.

We now show that $W_*$ has an ordered structure.
For each fixed $k$, averaging over the squares gives
\[
  \|W_n^{(k)}-W_*^{(k)}\|_1\le\|W_n-W_*\|_1\longrightarrow0.
\]
On this fixed partition, each block value of $W_n^{(k)}$
converges to the corresponding block value
of $W_*^{(k)}$, preserving the inequalities between block values.
Thus $W_*^{(k)}$ is also nondecreasing in each variable.
Letting $n\to\infty$ in the approximation bound $\|W_n-W_n^{(k)}\|_1\le2/k$
gives $\|W_*-W_*^{(k)}\|_1\le2/k$.
Thus $W_*^{(k)}\to W_*$ in $L^1([0,1]^2)$. Choose a subsequence
$(W_*^{(k_i)})_{i=1}^{\infty}$ that converges pointwise to $W_*$ on a
symmetric set $F\subseteq[0,1]^2$ of full measure.
For $(x,y),(x',y')\in F$ with $x\le x'$ and $y\le y'$,
since $W_*^{(k_i)}(x,y)\le W_*^{(k_i)}(x',y')$ for every $i$,
passing to the limit gives $W_*(x,y)\le W_*(x',y')$.
Hence the symmetric function
\[
  (x,y)\longmapsto
  \sup\bigl\{W_*(x',y'):(x',y')\in F,\ x'\le x,\ y'\le y\bigr\}
\]
with $\sup\emptyset:=0$ is nondecreasing in each variable and agrees with
$W_*$ on $F$. It differs from $W_*$ only on the null set $F^c$, so
$W_*$ has the required ordered structure.
\end{proof}

The uniform cost bound in the next lemma will control the left
difference quotients of $\Psi_\eps$ at the limit density. By
\Cref{lem:degenerate-slope}, this will exclude any degenerate minimizer there.

\begin{lemma}\label{lem:uniform-density-restoration}
Let $W_n\to W_*$ in $L^1([0,1]^2)$ be graphons, and suppose that $W_*$ satisfies the
constraint qualification.  There is $C<\infty$ such that, for all large $n$,
there is a graphon $V_n$ with
\[
  \mathbf C(V_n)=\mathbf C(W_*),\qquad
  |I_\eps(V_n)-I_\eps(W_n)|\le C|\mathbf C(W_n)-\mathbf C(W_*)|.
\]
\end{lemma}

\begin{proof}
Choose $\gamma,U_1,U_2$ satisfying \Cref{def:constraint-qualification}
at $W_*$, and restrict these directions to the interior region of $W_n$:
\[
  U_{i,n}:=U_i\mathbf1_{\{\gamma/2<W_n<1-\gamma/2\}},\qquad i=1,2.
\]
For $\xi=(\xi_1,\xi_2)^{\mathsf T}\in\RR^2$, define the bounded symmetric kernel
\[
  V_n(\xi):=W_n+\xi_1U_{1,n}+\xi_2U_{2,n}.
\]
Since $W_n\to W_*$ in $L^1([0,1]^2)$ and $U_i$ is supported where
$\gamma<W_*<1-\gamma$, we have $U_{i,n}\to U_i$ in $L^1([0,1]^2)$.
By~\eqref{eq:derivative-kernel-lipschitz}, we also have
$\mathcal D_{H,W_n}\to\mathcal D_{H,W_*}$ in $L^1([0,1]^2)$.
These convergences give
\[
  \Gamma_n:=\Gamma_{W_n}(U_{1,n},U_{2,n})
  \longrightarrow\Gamma_{W_*}(U_1,U_2).
\]
The limiting matrix is invertible, so there is $M<\infty$ such that
$\|\Gamma_n^{-1}\|_{\mathrm{op}}\le M$ for all large $n$.

Define the remainder in the linear approximation to the density map by
\begin{equation}\label{eq:remainder}
  R_n(\xi):=\mathbf C(V_n(\xi))-\mathbf C(W_n)-\Gamma_n\xi.
\end{equation}
The definition of $\Gamma_n$ gives $R_n(0)=0$ and $DR_n(0)=0$.
Since the densities are polynomial in $\xi$, the bounds $0\le W_n\le1$
and $\|U_{i,n}\|_\infty\le\|U_i\|_\infty$ imply that the second
derivatives of $R_n$ with respect to $\xi$ are uniformly bounded in $n$
on a fixed ball about the origin.
Hence, on that ball,
\[
  |R_n(\xi)|\le C|\xi|^2,\qquad \|DR_n(\xi)\|_{\mathrm{op}}\le C|\xi|,
\]
with $C$ independent of $n$.  Put
$\Delta_n:=\mathbf C(W_*)-\mathbf C(W_n)$, which tends to zero.
Using~\eqref{eq:remainder}, we can rewrite the constraint equations
$\mathbf C(V_n(\xi))=\mathbf C(W_*)$ as
\[
  \xi=\Gamma_n^{-1}\bigl(\Delta_n-R_n(\xi)\bigr).
\]
For all sufficiently large $n$ and $|\xi|\le2M|\Delta_n|$, we have
\begin{align*}
  \bigl|\Gamma_n^{-1}(\Delta_n-R_n(\xi))\bigr|
  &\le M|\Delta_n|+4CM^3|\Delta_n|^2\le2M|\Delta_n|,\\
  \|\Gamma_n^{-1}DR_n(\xi)\|_{\mathrm{op}}
  &\le2CM^2|\Delta_n|<\tfrac12.
\end{align*}
The first inequality shows that the map
$\xi\mapsto\Gamma_n^{-1}(\Delta_n-R_n(\xi))$ sends the closed ball
$|\xi|\le2M|\Delta_n|$ into itself. The second shows that its Lipschitz
constant on this ball is at most $1/2$. The contraction mapping theorem
therefore gives a fixed point $\xi_n=(\xi_{1,n},\xi_{2,n})^{\mathsf T}$ satisfying
\[
  |\xi_n|\le2M|\Delta_n|,\qquad
  \mathbf C(V_n(\xi_n))=\mathbf C(W_*).
\]
Set $V_n:=V_n(\xi_n)$.
For large $n$, the perturbation has $L^\infty([0,1]^2)$-norm at most
$\gamma/4$ and is supported in $\Omega_{\gamma/2}(W_n)$, so $V_n$ is a graphon.
On this set, both $W_n$ and $V_n$ take values in $[\gamma/4,1-\gamma/4]$.
Let
\[
  L:=\max_{z\in[\gamma/4,1-\gamma/4]}|J_\eps'(z)|<\infty.
\]
Since each $U_{i,n}$ vanishes outside $\Omega_{\gamma/2}(W_n)$, we have
$V_n=W_n$ outside this set. Applying the mean value theorem gives
\begin{align*}
  |I_\eps(V_n)-I_\eps(W_n)|
  &=\left|\int_{\Omega_{\gamma/2}(W_n)}
    \bigl(J_\eps(V_n)-J_\eps(W_n)\bigr)\right|\\
  &\le L\int_{\Omega_{\gamma/2}(W_n)}|V_n-W_n|\\
  &\le L\sum_{i=1}^2|\xi_{i,n}|
    \int_{\Omega_{\gamma/2}(W_n)}|U_{i,n}|\\
  &=L\sum_{i=1}^2|\xi_{i,n}|\,\|U_{i,n}\|_1\\
  &\le2ML\bigl(\|U_1\|_1+\|U_2\|_1\bigr)
    |\mathbf C(W_n)-\mathbf C(W_*)|.
\end{align*}
In the last step, we used $\|U_{i,n}\|_1\le\|U_i\|_1$ and
$|\xi_{i,n}|\le|\xi_n|\le2M|\Delta_n|$.
The resulting constant is independent of $n$.
\end{proof}

\begin{proof}[Proof of \Cref{prop:optimizer-nondegeneracy}]
Fix $\tau\in(\eps^m,\tau_{\max}(\eps))$.
\Cref{cor:almost-every-tau} gives a sequence $\tau_n\uparrow\tau$
at which every minimizer has ordered blocks. Choose minimizers $W_n$,
which therefore have an ordered structure. By \Cref{lem:ordered-compact},
after suitable relabelings we may pass to a subsequence converging in
$L^1([0,1]^2)$ to a graphon $W_*$ with an ordered structure.
By~\eqref{eq:density-l1-continuity}, $\mathbf C(W_*)=(\eps,\tau)$.
Feasibility of $W_*$, entropy continuity~\eqref{eq:entropy-l1-continuity},
and monotonicity from \Cref{lem:fixed-monotonicity} give
\[
  \Psi_\eps(\tau)\le I_\eps(W_*)
  =\lim_n I_\eps(W_n)=\lim_n\Psi_\eps(\tau_n)\le\Psi_\eps(\tau).
\]
Thus $W_*$ is a minimizer at $\tau$, and \Cref{lem:ordered-nondegenerate}
shows that it satisfies the constraint qualification.

We have obtained one nondegenerate minimizer at $\tau$; it remains to
exclude degenerate minimizers that might coexist with $W_*$.
Using the constraint qualification at $W_*$, \Cref{lem:uniform-density-restoration}
supplies graphons $V_n$ feasible at $\tau$ and a constant $C$ independent
of $n$ such that
\[
  \Psi_\eps(\tau)\le I_\eps(V_n)
  \le\Psi_\eps(\tau_n)+C(\tau-\tau_n).
\]
Thus the left difference quotients along this sequence are bounded:
\[
  \frac{\Psi_\eps(\tau)-\Psi_\eps(\tau_n)}{\tau-\tau_n}\le C.
\]
If any minimizer at $\tau$ were degenerate,
\Cref{lem:degenerate-slope} would force the left difference
quotients to diverge along every sequence approaching $\tau$ from
below.  This contradicts the bound along $(\tau_n)$.
Every minimizer at $\tau$ is therefore nondegenerate.
\end{proof}

\begin{proof}[Proof of \Cref{thm:fixed-densities}]
The result follows from
\Cref{prop:nondegenerate-multipodality,prop:optimizer-nondegeneracy}.
\end{proof}
